\section{Related work}\label{sec:related}

\paragraph{Conformal prediction for time series.} Conformal prediction \citep{vovk2005algorithmic} turns a point predictor into a prediction set with finite-sample coverage when the data are exchangeable; in the split variant, residuals on a held-out calibration set define a quantile and the interval is the point forecast plus or minus that quantile. Time series are not exchangeable and the guarantee does not hold in general. \citet{barber2023beyond} give a coverage bound beyond exchangeability, \citet{xu2021enbpi} an ensemble-based interval for dynamic time series, and \citet{zaffran2022agaci} study adaptive conformal prediction and the learning rate. A simple practical response is to compute the quantile from the most recent residuals only, which we call rolling conformal.

\paragraph{Adaptive methods.} Adaptive conformal inference (ACI; \citealp{gibbs2021aci}) updates the working miscoverage level online by a step size times the difference between the target miscoverage and the realised miss indicator, so that the long-run miss frequency approaches the target without assumptions on the data. FACI \citep{gibbs2024faci} aggregates several step sizes, conformal PID control \citep{angelopoulos2023pid} controls a quantile with proportional, integral and forecasting components, and SAOCP \citep{bhatnagar2023saocp} uses strongly adaptive online learning. The benchmark includes ACI, FACI and conformal PID; SAOCP is in the released \texttt{dracp} package and was run in a reproduction check only. These guarantees concern long-run miss frequency, not every country, horizon and window, and they are not stated for the situation in which the outcome of an $h$-step forecast is revealed only $h$ months later; this paper measures the effect of that delay (Section~\ref{sec:feedback}).

\paragraph{DRACP.} As described by its authors (Dynamic Regime-Aware Conformal Calibration; \citealp{oancea2026dracp}, version 2, read from the HTML version), DRACP keeps the split-conformal radius but replaces equal weights on past errors by a weighted quantile whose weights combine a temporal decay, a density-ratio correction for covariate shift, a Gaussian kernel with a floored effective sample size and a regime-similarity term from a Gaussian mixture model, and an online controller adjusts the level over several learning rates. The authors evaluate it on 48 series (28 euro-area and EU HICP inflation series, 13 US macroeconomic and energy series, 7 daily financial series), mainly one step ahead with an additional multi-horizon evaluation at $h=3,6,12$; they report mean coverage over 48 series and 20 seeds of 0.890 for DRACP, 0.884 for ACI, 0.883 for FACI, 0.868 for SAOCP, 0.841 for rolling, 0.825 for conformal PID and 0.781 for split at nominal 0.90, and state that the advantage is specific to one-step-ahead forecasting and depends on the point forecaster. This description was read from the authors' paper by the project author and is to be checked against it before submission. Our setting differs in three respects: the point model is refitted on an expanding window, the target is the cumulative $h$-month log growth, and the economies are commodity-dependent. The study is a transfer of a published method to a different panel and draws no conclusion about the method beyond the limitation its authors state.

\paragraph{Inflation uncertainty and pretrained models.} Intervals for inflation are commonly derived from the forecast model's own error distribution, and we compare conformal calibration with the native ARIMA and ETS intervals in a secondary, post hoc analysis (Section~\ref{sec:native}). We ran no systematic literature search on forecast uncertainty in emerging markets, make no claim about that literature and cite no study of the calibration of inflation intervals there, because none was verified. Chronos \citep{ansari2024chronos} is a family of pretrained time-series models that can be used zero-shot; we use the smallest Chronos-Bolt model as a point forecaster with conformal intervals on top (Section~\ref{sec:forecasters}), make no claim about foundation models beyond this use, and mark its bibliographic entry as needing verification.
